%% apendices/exemplo_numerico.tex
%% Apêndice A — Exemplo numérico da ordem de achatamento da camada densa

\chapter{Exemplo numérico da ordem de achatamento da camada densa}
\label{ap:exemplo-numerico}

Este apêndice complementa a Seção~\ref{sec:camada-densa}, rastreando,
elemento a elemento, a correspondência entre a matriz de pesos de uma camada
\texttt{Linear} do PyTorch e a sequência linear sobre a qual a SampEn$_{1D}$ é
calculada. Usa-se uma camada mínima, $2 \times 3$, primeiro na ordem por
neurônio (Equação~\ref{eq:ordem-neuronio}) e em seguida na ordem por entrada
(Equação~\ref{eq:ordem-entrada}).

% ===========================================================================
\section{Camada Linear: matriz \texorpdfstring{$2 \times 3$}{2x3}}
\label{sec:ap-linear}

Considere uma camada \texttt{nn.Linear(in\_features=3, out\_features=2)}
cujo tensor \texttt{weight} de \emph{shape} $[\mathrm{Out}, \mathrm{In}] = [2, 3]$
recebe valores inteiros para facilitar o rastreamento:
\begin{equation}
  W_{\text{lin}} \;=\;
  \begin{pmatrix}
    w_{0,0} & w_{0,1} & w_{0,2} \\
    w_{1,0} & w_{1,1} & w_{1,2}
  \end{pmatrix}
  \;=\;
  \begin{pmatrix}
    1 & 2 & 3 \\
    4 & 5 & 6
  \end{pmatrix}.
\end{equation}

Aplicando \texttt{W\_lin.flatten()} obtemos o vetor

\begin{equation}
  W_{\text{lin}}.\texttt{flatten()} \;=\; [\,1,\;2,\;3,\;4,\;5,\;6\,].
\end{equation}

A Tabela~\ref{tab:ap-linear} mostra a correspondência elemento a elemento.
Note que a leitura é feita \textbf{linha por linha}: primeiro saem todos os
três pesos de entrada do neurônio de saída $\mathrm{Out}=0$, depois os três
pesos do neurônio $\mathrm{Out}=1$. Cada linha contígua do vetor
corresponde a \emph{um único neurônio de saída}.

\begin{table}[H]
\centering
\caption{Mapeamento elemento a elemento da camada Linear $2 \times 3$.}
\label{tab:ap-linear}
\small
\begin{tabular}{cccccc}
\toprule
\textbf{Posição no vetor} & 0 & 1 & 2 & 3 & 4 \\
                          & ($\to$) 5 & & & & \\
\midrule
\textbf{Elemento}      & $w_{0,0}$ & $w_{0,1}$ & $w_{0,2}$ & $w_{1,0}$ & $w_{1,1}$ \\
                       & $w_{1,2}$ & & & & \\
\midrule
\textbf{Valor}         & 1 & 2 & 3 & 4 & 5 \\
                       & 6 & & & & \\
\midrule
\textbf{Pertence a}    & \multicolumn{3}{c}{\textbf{Saída $\mathrm{Out}=0$}} & \multicolumn{2}{c}{\textbf{Saída $\mathrm{Out}=1$}} \\
\bottomrule
\end{tabular}
\end{table}

% ===========================================================================
\section{Contraste: ordem por entrada}
\label{sec:ap-transposicao}

Transpondo a mesma matriz antes do achatamento, percorrem-se primeiro todos os
pesos que partem da entrada $\mathrm{In}=0$, depois os da entrada
$\mathrm{In}=1$, e assim por diante:

\begin{equation}
  W_{\text{lin}}^{\top}.\texttt{flatten()} \;=\; [\,1,\;4,\;2,\;5,\;3,\;6\,].
\end{equation}

Os dois vetores contêm exatamente os mesmos seis valores, e portanto o mesmo
histograma: a complexidade LMC é idêntica nas duas ordens. O que muda é a
vizinhança de cada valor, à qual a SampEn$_{1D}$ é sensível.

% ===========================================================================
\section{Reprodutibilidade}
\label{sec:ap-reproducao}

Os vetores apresentados neste apêndice podem ser reproduzidos diretamente em
PyTorch:

\begin{verbatim}
    import torch

    # Linear (Out=2, In=3)
    W_lin = torch.tensor([[1., 2., 3.],
                          [4., 5., 6.]])
    print(W_lin.flatten().tolist())      # ordem por neuronio
    # [1.0, 2.0, 3.0, 4.0, 5.0, 6.0]
    print(W_lin.t().flatten().tolist())  # ordem por entrada
    # [1.0, 4.0, 2.0, 5.0, 3.0, 6.0]
\end{verbatim}
